\documentclass[aspectratio=169,10pt,spanish]{beamer}

\input{../../../_preambulo_beamer.tex}
\input{../../../_comandos_beamer.tex}

\selectlanguage{spanish}

\title{MA1001B - Modelación Estadística para la Toma de Decisiones}
\subtitle{Lección 02: Diseños de muestreo probabilístico}
\author{}
\institute{Tecnol\'ogico de Monterrey}
\date{}

\begin{document}

\begin{frame}[plain]
  \vspace{1.2cm}
  {\Large\bfseries MA1001B - Modelación Estadística para la Toma de Decisiones\par}
  \vspace{0.7cm}
  {\large Lección 02: Diseños de muestreo probabilístico\par}
  \vspace{0.7cm}
  {\color{TecRojo}\rule{\textwidth}{1.2pt}}
  \vspace{0.8cm}

  Facultad MA1001B\\[0.3cm]
  Periodo\\[0.3cm]
  Tecnol\'ogico de Monterrey
\end{frame}

\begin{frame}{La pregunta de diseño muestral}
  \begin{alertblock}{Pregunta guía}
    ¿Qué unidades se seleccionarán al azar y qué partes de la población cubrirá
    la muestra resultante?
  \end{alertblock}

  \vspace{0.5em}
  Al final de esta lección, usted puede:
  \begin{itemize}
    \item conectar una muestra probabilística con un marco muestral;
    \item implementar muestreo aleatorio simple y estratificado proporcional;
    \item identificar la unidad seleccionada en el muestreo por conglomerados
      de una etapa;
    \item comparar diseños sin declarar un ganador universal.
  \end{itemize}

  \vfill
  \scriptsize Todos los registros son sintéticos. No se usan datos reales de empresa.
\end{frame}

\begin{frame}{Tres diseños de muestreo probabilístico}
  \small
  \begin{center}
    \begin{tabular}{p{0.20\textwidth}p{0.27\textwidth}p{0.39\textwidth}}
      \toprule
      \textbf{Diseño} & \textbf{Unidad aleatorizada} & \textbf{Construcción} \\
      \midrule
      Aleatoria simple & Pedido individual & Seleccionar $n$ pedidos de un marco completo \\
      Estratificado & Pedido dentro de una región & Partir en estratos; muestrear de cada estrato \\
      Conglomerados de una etapa & Instalación & Seleccionar instalaciones; incluir todos sus pedidos \\
      \bottomrule
    \end{tabular}
  \end{center}

  \vspace{0.7em}
  \begin{exampleblock}{Marco sintético}
    $12{,}000$ pedidos $=$ 24 instalaciones $\times$ 500 pedidos. Cuatro regiones
    definen estratos; las instalaciones definen conglomerados.
  \end{exampleblock}

  \textbf{Diseño probabilístico:} la selección se rige por un proceso de azar
  conocido, no por qué registros resultan más fáciles de alcanzar.
\end{frame}

\begin{frame}[fragile]{Paso 1: Muestreo aleatorio simple}
  \begin{columns}[T]
    \begin{column}{0.42\textwidth}
      \small
      \begin{lstlisting}
sample = population.sample(
    n=240,
    replace=False,
    random_state=42,
)

inclusion_chance = n / N
      \end{lstlisting}

      \begin{block}{Salida verificada}
        $N=12{,}000$, $n=240$\\
        Probabilidad de inclusión $=2.00\%$\\
        $\mu=47.67$, $\bar{x}=47.17$ min\\
        Error absoluto $=0.50$ min
      \end{block}
    \end{column}
    \begin{column}{0.58\textwidth}
      \centering
      \includegraphics[width=\textwidth]{../figuras/disenos_muestreo_01_aleatorio_simple.png}
      \vspace{0.15em}
      \scriptsize La probabilidad igual de selección no fuerza un balance
      regional exacto.
    \end{column}
  \end{columns}
\end{frame}

\begin{frame}[fragile]{Paso 2: Muestreo estratificado proporcional}
  \begin{columns}[T]
    \begin{column}{0.40\textwidth}
      \small
      \textbf{Regla de asignación}
      \[
        n_h = n\frac{N_h}{N}
      \]
      Aquí $N_h=3{,}000$ para cada región, así que $n_h=60$.

      \vspace{0.4em}
      \begin{lstlisting}
for region in regions:
    stratum = population[
        population.region == region
    ]
    sample 60 orders
      \end{lstlisting}

      Error del MAS: 0.50 min\\
      Error estratificado: 0.28 min
    \end{column}
    \begin{column}{0.60\textwidth}
      \centering
      \includegraphics[width=\textwidth]{../figuras/disenos_muestreo_02_estratificado.png}
      \vspace{0.15em}
      \scriptsize La estratificación garantiza representación, no el menor
      error en cada muestra posible.
    \end{column}
  \end{columns}
\end{frame}

\begin{frame}{Paso 3: Muestreo por conglomerados de una etapa}
  \begin{columns}[T]
    \begin{column}{0.40\textwidth}
      \small
      Instalaciones seleccionadas al azar:
      \[
        \{\mathrm{F02},\mathrm{F11},\mathrm{F16},\mathrm{F18}\}
      \]
      Incluir los 500 pedidos de cada instalación seleccionada.

      \vspace{0.5em}
      \begin{tabular}{lr}
        \toprule
        Pedidos incluidos & 2{,}000 \\
        Probabilidad de inclusión & 16.67\% \\
        Regiones presentes & 3 de 4 \\
        Estimación de la media & 48.76 min \\
        Error absoluto & 1.09 min \\
        \bottomrule
      \end{tabular}

      \vspace{0.5em}
      Unidades cercanas pueden compartir un efecto de instalación; una muestra
      grande de conglomerados puede aportar menos información independiente de
      lo que sugiere su número de filas.
    \end{column}
    \begin{column}{0.60\textwidth}
      \centering
      \includegraphics[width=\textwidth]{../figuras/disenos_muestreo_03_conglomerados.png}
      \vspace{0.15em}
      \scriptsize Las barras rosa marcan las instalaciones seleccionadas en
      esta muestra con semilla.
    \end{column}
  \end{columns}
\end{frame}

\begin{frame}{Elegir un diseño}
  \small
  \begin{center}
    \begin{tabular}{p{0.18\textwidth}p{0.23\textwidth}p{0.22\textwidth}p{0.22\textwidth}}
      \toprule
      \textbf{Diseño} & \textbf{Marco necesario} & \textbf{Cobertura} & \textbf{Compensación operativa} \\
      \midrule
      Aleatoria simple & Lista completa de pedidos & Aleatoria, no fija por región & Las unidades pueden estar dispersas \\
      Estratificado & Lista de pedidos más región & Cada región muestreada & Requiere etiquetas de región \\
      Conglomerados & Lista de instalaciones & Solo instalaciones seleccionadas & Menos sitios que visitar \\
      \bottomrule
    \end{tabular}
  \end{center}

  \vspace{0.8em}
  \begin{alertblock}{Interprete la comparación correctamente}
    Los errores observados---0.50, 0.28 y 1.09 minutos---pertenecen a una
    realización sintética con semilla. Ilustran consecuencias de la estructura
    poblacional; no ordenan los tres diseños para todo contexto.
  \end{alertblock}
\end{frame}

\begin{frame}{Verificación de conceptos}
  \begin{enumerate}
    \item En el MAS, ¿cuál es la probabilidad de inclusión de un pedido y qué
      marco se requiere para implementar el diseño?
    \item ¿Por qué la muestra estratificada contiene cada región mientras que
      el MAS no garantiza conteos regionales exactos?
    \item En el diseño por conglomerados, ¿qué se selecciona al azar y qué
      pedidos entran después a la muestra?
    \item ¿Por qué esta comparación con semilla no prueba que la estratificación
      sea siempre el diseño más exacto?
  \end{enumerate}

  \vspace{0.5em}
  \begin{block}{Conexión con la Lección 03}
    Una vez recolectada una muestra defendible, los resúmenes numéricos
    describen su centro y su variabilidad.
  \end{block}

  \vfill
  \scriptsize Fuente: OpenStax, \textit{Introductory Business Statistics 2e},
  Capítulo 1, Sección 1.2. CC BY-NC-SA 4.0.
\end{frame}

\appendix



\end{document}
